% Options for packages loaded elsewhere
\PassOptionsToPackage{unicode}{hyperref}
\PassOptionsToPackage{hyphens}{url}
\PassOptionsToPackage{dvipsnames,svgnames,x11names}{xcolor}
\documentclass[
  a4paper,
]{article}
\usepackage{xcolor}
\usepackage[margin=25mm]{geometry}
\usepackage{amsmath,amssymb}
\setcounter{secnumdepth}{-\maxdimen} % remove section numbering
\usepackage{iftex}
\ifPDFTeX
  \usepackage[T1]{fontenc}
  \usepackage[utf8]{inputenc}
  \usepackage{textcomp} % provide euro and other symbols
\else % if luatex or xetex
  \usepackage{unicode-math} % this also loads fontspec
  \defaultfontfeatures{Scale=MatchLowercase}
  \defaultfontfeatures[\rmfamily]{Ligatures=TeX,Scale=1}
\fi
\usepackage{lmodern}
\ifPDFTeX\else
  % xetex/luatex font selection
  \setmainfont[]{DejaVu Serif}
  \setsansfont[]{DejaVu Sans}
  \setmonofont[]{DejaVu Sans Mono}
\fi
% Use upquote if available, for straight quotes in verbatim environments
\IfFileExists{upquote.sty}{\usepackage{upquote}}{}
\IfFileExists{microtype.sty}{% use microtype if available
  \usepackage[]{microtype}
  \UseMicrotypeSet[protrusion]{basicmath} % disable protrusion for tt fonts
}{}
\makeatletter
\@ifundefined{KOMAClassName}{% if non-KOMA class
  \IfFileExists{parskip.sty}{%
    \usepackage{parskip}
  }{% else
    \setlength{\parindent}{0pt}
    \setlength{\parskip}{6pt plus 2pt minus 1pt}}
}{% if KOMA class
  \KOMAoptions{parskip=half}}
\makeatother
\usepackage{longtable,booktabs,array}
\newcounter{none} % for unnumbered tables
\usepackage{calc} % for calculating minipage widths
% Correct order of tables after \paragraph or \subparagraph
\usepackage{etoolbox}
\makeatletter
\patchcmd\longtable{\par}{\if@noskipsec\mbox{}\fi\par}{}{}
\makeatother
% Allow footnotes in longtable head/foot
\IfFileExists{footnotehyper.sty}{\usepackage{footnotehyper}}{\usepackage{footnote}}
\makesavenoteenv{longtable}
\usepackage{graphicx}
\makeatletter
\newsavebox\pandoc@box
\newcommand*\pandocbounded[1]{% scales image to fit in text height/width
  \sbox\pandoc@box{#1}%
  \Gscale@div\@tempa{\textheight}{\dimexpr\ht\pandoc@box+\dp\pandoc@box\relax}%
  \Gscale@div\@tempb{\linewidth}{\wd\pandoc@box}%
  \ifdim\@tempb\p@<\@tempa\p@\let\@tempa\@tempb\fi% select the smaller of both
  \ifdim\@tempa\p@<\p@\scalebox{\@tempa}{\usebox\pandoc@box}%
  \else\usebox{\pandoc@box}%
  \fi%
}
% Set default figure placement to htbp
\def\fps@figure{htbp}
\makeatother
\setlength{\emergencystretch}{3em} % prevent overfull lines
\providecommand{\tightlist}{%
  \setlength{\itemsep}{0pt}\setlength{\parskip}{0pt}}
\usepackage{luatexja}
\usepackage{xurl}
\usepackage{bookmark}
\IfFileExists{xurl.sty}{\usepackage{xurl}}{} % add URL line breaks if available
\urlstyle{same}
\hypersetup{
  pdftitle={Did We Understand Complex Numbers Correctly?},
  pdfauthor={Noriaki Kihara (WF System Co., Ltd.) ORCID 0009-0004-6753-4020},
  colorlinks=true,
  linkcolor={Maroon},
  filecolor={Maroon},
  citecolor={Blue},
  urlcolor={Blue},
  pdfcreator={LaTeX via pandoc}}

\title{Did We Understand Complex Numbers Correctly?}
\usepackage{etoolbox}
\makeatletter
\providecommand{\subtitle}[1]{% add subtitle to \maketitle
  \apptocmd{\@title}{\par {\large #1 \par}}{}{}
}
\makeatother
\subtitle{― Where the Amplitude Lives, the Imaginary Reading, and
Attribution across \(\mathbb{C}\) → \(\mathbb{H}\) → \(\mathbb{O}\) ―}
\author{Noriaki Kihara (WF System Co., Ltd.) ORCID 0009-0004-6753-4020}
\date{October 23, 2026 Version DOI 10.5281/zenodo.23585437}

\ltjsetparameter{jacharrange={-2,-3}}
\begin{document}
\maketitle

\textbf{Author:} Noriaki Kihara (WF System Co., Ltd.)　ORCID
0009-0004-6753-4020\\
\textbf{Date:} 2026-10-23　　\textbf{Version:} v1\\
\textbf{Version DOI:} 10.5281/zenodo.23585437　　\textbf{Concept DOI:}
10.5281/zenodo.23585436\\
\textbf{Same folder:} figures and identity checks in \texttt{図表/\allowbreak{}}
(\texttt{make\_\allowbreak{}figures.py}, \texttt{verify\_\allowbreak{}identities.py},
\texttt{results/\allowbreak{}checks.txt}; English-label figures \texttt{fig*\_\allowbreak{}en}),
citation tables \texttt{自己引用の対応表\_\allowbreak{}20261023.md} and
\texttt{外部引用の対応表\_\allowbreak{}20261023.md}, list of issues
\texttt{課題の整理\_\allowbreak{}20261023.md}\\
\textbf{Note:} The Japanese version
(\texttt{complex\_\allowbreak{}revisited\_\allowbreak{}amplitude\_\allowbreak{}imaginary\_\allowbreak{}ja\_\allowbreak{}v1.md}) is the
authoritative text; this English version is a translation.

\begin{center}\rule{0.5\linewidth}{0.5pt}\end{center}

\subsection{Abstract}\label{abstract}

The author's research programme aims to read out spacetime and particles
from the relations alone of an anonymous set of complex numbers \{zₙ\},
and has used the structured number systems ℝ → \(\mathbb{C}\) →
\(\mathbb{H}\) → \(\mathbb{O}\) in turn. This paper pauses midway and
checks whether a symmetry already present at the stage of complex
numbers has been overlooked. Without this check, a property found with
quaternions or octonions is misattributed as first appearing at that
stage. One rule is set for the check. \textbf{A property may be said to
belong to stage L only when it has been shown that it cannot be stated
with the algebra of stage L − 1 alone.}

The check starts from one naive question. When a state lies on a circle
(or a spherical shell), where does the amplitude live? The main contents
are as follows.

\begin{enumerate}
\def\labelenumi{\arabic{enumi}.}
\tightlist
\item
  On a circle in the plane the amplitude \textbar z\textbar{} and the
  radius of curvature R are the same number, so one does not notice that
  they are different things. The radius of curvature R is invisible
  under translation along a geodesic and can be read only by a closed
  loop that sweeps an area (holonomy). If the amplitude is placed on the
  radial direction, a change of amplitude becomes a change of curvature.
\item
  The amplitude therefore lives in an unreadable direction orthogonal
  both to the circle and to R. This is the origin of the Q axis of the
  author's spacetime equation x² ＋ y² ＋ z² ＋ (it)² ＋ (iR)² ＋ (iQ)²
  ＝ 0. The three tangential directions are space, the three normal
  directions (t, R, Q) are unreadable, and unreadable directions are
  written with imaginary units.
\item
  The amplitude in the unreadable direction is read as imaginary. For
  the logarithm w ＝ log u ＝ φ ＋ iθ of the relation u (φ ＝
  log-amplitude, θ ＝ phase), \textbar cosh(w/2)\textbar² ＝ (cosh φ ＋
  cos θ)/2 and \textbar sinh(w/2)\textbar² ＝ (cosh φ − cos θ)/2, and
  their sum is not 1 but cosh φ ＝ γ (the Lorentz factor). Only on the
  section φ ＝ 0 does cos²(θ/2) ＋ sin²(θ/2) ＝ 1 hold, and this is the
  form of the Born rule. The normalization that makes it look like a
  probability is division by γ, and the one dimension discarded there is
  the rapidity.
\item
  The denominator of the coupling coefficient k ＝ (1 − u)/(1 ＋ u) ＝
  −tanh(w/2) of the previous paper is \textbar1 ＋ u\textbar² ＝
  2e\^{}φ(cosh φ ＋ cos θ) ＝ 4e\^{}φ \textbar cosh(w/2)\textbar², the
  same division as the Born-rule normalization. \textbar k\textbar{} ＝
  1 holds when cos θ ＝ 0, independently of the amplitude mismatch φ.
  The point that the previous paper identified as a relation on the
  light cone is the self-dual point of coupling and distance.
\item
  Amplitude and phase are conjugate and dual, and the duality is the 90°
  rotation of the exponent plane by multiplication with i. Number--phase
  conjugacy, time and frequency, Wick rotation and the Laplace variable
  are appearances of the same rotation. Which axis is read as the
  imaginary part, and where one divides to normalize, cannot be derived
  from dynamics (self-reference) and is placed as an axiom. The Born
  rule and the invariance of the speed of light are two sections of this
  single axiom.
\item
  In complex numbers an amplitude read as imaginary cannot be
  distinguished from the phase (there is only one unit). In quaternions
  the amplitude can be given a separate unit m ⊥ n, so amplitude and
  phase are distinguished. Moreover they do not commute: the commutator
  {[}e\^{}\{ma\}, e\^{}\{nθ\}{]} ＝ 2 sin a sin θ (m × n) generates a
  third direction. The three fill the three-dimensional imaginary part,
  that is, space. This is the root of the previous paper's k being a
  three-dimensional vector, and k is the sign-reversed modified
  Rodrigues parameter.
\item
  In octonions, when the units of amplitude are two (m for R and l for
  Q) and form a triple with the phase unit n, the associator first
  becomes non-zero, and an R-type amplitude and a Q-type amplitude
  become algebraically different things. With the spacetime-side
  quaternion a ＝ t ＋ xi ＋ yj ＋ zk and the internal-side quaternion b
  ＝ R ＋ Q₁i ＋ Q₂j ＋ Q₃k, the same 1 ＋ 3 split repeats. Reading B:
  mass ＝ spacetime curvature R (the scalar part of the internal side),
  charge ＝ internal curvature Q (the vector part of the internal side).
  Amplitude ＝ \textbar b\textbar. Hypothesis 1 of the previous paper
  (mass ∝ norm) is refined by this reading.
\item
  All of this is collected in an attribution ledger: what \(\mathbb{C}\)
  already has (duality, imaginary amplitude and the light cone, cos² and
  γ, the self-dual point, U\^{}n ＝ I and m/n, the division trap,
  amplitude ≠ radius of curvature), what \(\mathbb{H}\) adds (the
  distinction of amplitude and phase, the third direction, the light
  cone with an orientation, memory of position, the possibility of
  condensates), and what \(\mathbb{O}\) adds (the associator, the
  distinction of mass-type and charge-type, the split form and
  \textbar a\textbar{} ＝ \textbar b\textbar, the order of binding).
\end{enumerate}

All identities used in this paper were checked numerically (11 items,
\texttt{図表/\allowbreak{}results/\allowbreak{}checks.txt}). The classification of claims and the
form of falsification are in §9.

\textbf{Keywords:} complex numbers, quaternions, octonions, amplitude
and phase, conjugacy, duality, imaginary reading, rapidity, Born rule,
Lorentz factor, modified Rodrigues parameters, associator, anonymity,
partial projection, attribution

\begin{center}\rule{0.5\linewidth}{0.5pt}\end{center}

\subsection{1. Research motivation}\label{research-motivation}

\subsubsection{1.1 The starting question}\label{the-starting-question}

General relativity and the Standard Model cannot be unified, yet when
approximated as locally flat both hold to extremely high precision. Why?
Is there a way to explain this without introducing arbitrary parameters?
This is the starting point of the author's research programme, and the
series of papers up to the previous paper {[}S30{]} has approached this
question by the method of reading out spacetime and particles from the
relations alone of an anonymous set of complex numbers \{zₙ\} {[}S1{]}
{[}S3{]} {[}S18{]} {[}S27{]} {[}S30{]}.

\subsubsection{1.2 Why return to complex
numbers}\label{why-return-to-complex-numbers}

The research has used the structured number systems ℝ → \(\mathbb{C}\) →
\(\mathbb{H}\) → \(\mathbb{O}\) in turn. Starting from the geometric
formulation of four-dimensional space by central projection {[}S29{]},
it considered the closure Σzₙ² ＝ 0 with complex numbers {[}S4{]}
{[}S28{]}, built a quaternion from two complex numbers to read out (1 ＋
3) {[}S27{]}, proposed inflation as crystallization from the relations
of quaternions {[}S30{]}, and wrote the spacetime equation with split
octonions {[}S30 §11{]}.

Here a doubt arose. Had the symmetry already present at the stage of
complex numbers been understood correctly? If a symmetry present from
the stage of complex numbers has been overlooked, a property found with
quaternions or octonions is misattributed as one that first appeared at
that stage. If the attribution is wrong, the answers to the questions
``why quaternions?'' and ``why octonions?'' are wrong too.

This paper therefore goes back and checks before going on. One rule is
set for the check.

\begin{quote}
\textbf{Decision rule.} A property may be said to belong to stage L
(\(\mathbb{C}\), \(\mathbb{H}\), \(\mathbb{O}\)) only when it has been
shown that it cannot be stated with the algebra of stage L − 1 alone.
For each claim, first try to write it one stage down; if it can be
written there, attribute it to the lower stage.
\end{quote}

\subsubsection{1.3 One naive question}\label{one-naive-question}

The entrance to the check is a naive question. When a complex number z
＝ a ＋ ib is read as a circle in the complex plane, the amplitude is
simply the real radius. We have come this far without finding that
strange. But if the state is taken to lie on a circle of radius R (for
quaternions, the four-dimensional spherical shell S³), then R is the
radius of curvature of that circle and should be a quantity different
from the amplitude. Where does the amplitude live?

This paper proceeds along this question. §3 decides where the amplitude
lives; §4 examines what comes out at the stage of complex numbers when
the amplitude is read as imaginary; §5 examines what quaternions add and
§6 what octonions add, checking with the decision rule as we go. §7
gives the attribution ledger, §8 the comparison with prior work, §9 the
classification of claims and the form of falsification, and §10 the
conclusion.

In each section, \textbf{what is mathematically trivial (identities)}
and \textbf{what is chosen as a physical reading (which section to read,
where to divide)} are written separately.

\begin{center}\rule{0.5\linewidth}{0.5pt}\end{center}

\subsection{2. Conventions}\label{conventions}

\subsubsection{2.1 Inherited from the previous
paper}\label{inherited-from-the-previous-paper}

The axioms and definitions are those of the previous paper {[}S30{]} and
the catalogue {[}S3{]}, used as they are. Anonymity (elements have no
names; there is no special origin, coordinate, distance, sign or kind)
{[}S1{]}; self-consistency (the system takes itself as its own
reference) {[}S2{]}; distinguishability (a readout cannot tell to which
layer a value belongs) {[}S3{]}; partial projection {[}S4{]} {[}S15{]}.
The linear closure Σq ＝ 0 and the true-square closure Σq² ＝ 0 are
treated as results, not conditions {[}S2{]} {[}S28{]}.

The quaternion notation (q ＝ s ＋ v, polar form q ＝
\textbar q\textbar{} e\^{}\{nθ\}, conjugate q̄, true square q² ＝ (s² −
\textbar v\textbar²) ＋ 2sv), the relative phase u ＝ q₁⁻¹q₂, the
coupling coefficient k ＝ (q₁ ＋ q₂)⁻¹(q₁ − q₂) and the distance x ＝
k⁻¹ are as in the previous papers {[}S27{]} {[}S30{]}.

\subsubsection{2.2 Notation, and the declarations of division, origin
and
observer}\label{notation-and-the-declarations-of-division-origin-and-observer}

This system is a system of products. Division is an implicit
normalization, so it is done once only, at the end of a readout, and its
denominator is declared each time. Addition inside a product moves the
origin, so the origin is declared whenever addition is used. Every value
carries an observer, that is, whose value it is {[}S5{]} {[}S6{]}. The
declarations for the quantities used in this paper are as follows.

{\def\LTcaptype{none} % do not increment counter
\begin{longtable}[]{@{}
  >{\raggedright\arraybackslash}p{(\linewidth - 8\tabcolsep) * \real{0.2000}}
  >{\raggedright\arraybackslash}p{(\linewidth - 8\tabcolsep) * \real{0.2000}}
  >{\raggedright\arraybackslash}p{(\linewidth - 8\tabcolsep) * \real{0.2000}}
  >{\raggedright\arraybackslash}p{(\linewidth - 8\tabcolsep) * \real{0.2000}}
  >{\raggedright\arraybackslash}p{(\linewidth - 8\tabcolsep) * \real{0.2000}}@{}}
\toprule\noalign{}
\begin{minipage}[b]{\linewidth}\raggedright
Quantity
\end{minipage} & \begin{minipage}[b]{\linewidth}\raggedright
Definition
\end{minipage} & \begin{minipage}[b]{\linewidth}\raggedright
Denominator
\end{minipage} & \begin{minipage}[b]{\linewidth}\raggedright
Origin
\end{minipage} & \begin{minipage}[b]{\linewidth}\raggedright
Observer
\end{minipage} \\
\midrule\noalign{}
\endhead
\bottomrule\noalign{}
\endlastfoot
Relative phase u & u ＝ q₁⁻¹q₂ (left division) & q₁ & --- & q₁ \\
Logarithm w & w ＝ log u ＝ φ ＋ iθ (complex case). φ ＝ log-amplitude
＝ rapidity, θ ＝ phase & --- & --- & q₁ \\
Complex angle ω & ω ＝ θ − iφ, u ＝ e\^{}\{iω\}, w ＝ iω & --- & --- &
q₁ \\
Coupling coefficient k & k ＝ (1 ＋ u)⁻¹(1 − u) ＝ (q₁ ＋ q₂)⁻¹(q₁ − q₂)
& the composite 1 ＋ u (＝ q₁ ＋ q₂) & the origin O of sum and
difference (the centroid fixed by Σq ＝ 0) & q₁ \\
Distance D & D ＝ 1/\textbar k\textbar{} & k & as above & q₁ \\
Quaternion exponentials & e\^{}\{ma\}, e\^{}\{nθ\} (m, n unit pure
imaginary, m ⊥ n). The compositions e\^{}\{ma\} e\^{}\{nθ\} and
e\^{}\{nθ\} e\^{}\{ma\} are distinguished, and both are distinguished
from e\^{}\{ma＋nθ\} (§5.2) & --- & --- & declared by the order of
multiplication \\
\end{longtable}
}

Written in w, cosh(w/2) ＝ cos(ω/2) and sinh(w/2) ＝ i sin(ω/2). This
paper mainly uses w ＝ log u; the complex angle ω is used only to show
the correspondence.

\subsubsection{2.3 Notation for
attribution}\label{notation-for-attribution}

At the end of each claim, the attribution by the decision rule is marked
【\(\mathbb{C}\)】, 【\(\mathbb{C}\) × ℝ】, 【\(\mathbb{H}\)】 or
【\(\mathbb{O}\)】. 【\(\mathbb{C}\) × ℝ】 means something that becomes
visible once a single unreadable axis is added to the complex numbers
(no quaternions are needed).

\begin{center}\rule{0.5\linewidth}{0.5pt}\end{center}

\subsection{3. Where does the amplitude
live?}\label{where-does-the-amplitude-live}

\subsubsection{3.1 Invisible on a circle in the
plane}\label{invisible-on-a-circle-in-the-plane}

Put the complex number z ＝ R e\^{}\{iθ\} on a circle of radius R; the
amplitude is \textbar z\textbar{} ＝ R. Because the amplitude and the
radius of the circle are the same number, one does not notice that they
are different quantities. Reading the amplitude as ``the radial
component'' produces no contradiction (Fig. 2, left).

\subsubsection{3.2 The radius of curvature R can be read only by the
area of a closed
loop}\label{the-radius-of-curvature-r-can-be-read-only-by-the-area-of-a-closed-loop}

Suppose the state lies on a spherical shell S³ of radius R. The local
coordinates are three-dimensional, and R itself is not visible.
Translation along a geodesic (a great circle) shows zero effect of
curvature; only a two-dimensional motion that sweeps an area, that is,
going once round a closed loop, makes it observable as rotation ＝ swept
area / R² (holonomy, Gauss--Bonnet {[}X3{]}, geometric phase {[}X2{]}).
In the author's system the c ＝ 1 area sweep is the clock {[}S7{]}
{[}S8{]}, and curvature oscillations are read out from the area
cross-term of complex rotating waves {[}S9{]}. The distortion of the
geodesic unit cell in positively curved space also appears as an
evaluation of areas and volumes {[}S10{]}. §5 of the previous paper
{[}S30{]} stated that a localized observer cannot see a uniform
curvature and sees only tides. This is the same thing in other words.

\subsubsection{3.3 Why the amplitude cannot be placed on the radial
direction}\label{why-the-amplitude-cannot-be-placed-on-the-radial-direction}

If the amplitude is placed radially, the amplitude is a quantity on the
same axis as R. But R is a bending that can be read only by closed
loops, and a change of amplitude would then be a change of curvature. If
Hypothesis 1 of the previous paper {[}S30{]} (mass ∝
\textbar q\textbar²) is read radially, two states of different mass lie
on shells of different radius of curvature and cannot share a common
curvature scale inside the same condensate. The limits of the radius of
curvature {[}S11{]}, the degeneracy at R ＝ 0 {[}S12{]} and the
self-consistent determination of R {[}S13{]} all presuppose that R is a
quantity determined independently of the amplitude. At the time of the
previous paper it was already understood that the radius of curvature R
and the amplitude \textbar q\textbar{} of a state must not be identified
(the caution in {[}S30{]} §9), but where the amplitude lives had not
been decided.

{[}S14{]}, which discussed the stage for amplitudes and the choice of
the time direction, and the readout hierarchy of {[}S15{]}, which showed
that position had never been read, were also asking where in the
hierarchy of readouts the amplitude is placed.

\subsubsection{3.4 The direction orthogonal both to the circle and to R:
the origin of the Q
axis}\label{the-direction-orthogonal-both-to-the-circle-and-to-r-the-origin-of-the-q-axis}

If the amplitude cannot live radially, the remaining place is the
direction orthogonal both to the circle and to R, an axis standing out
of the plane of the circle. The circle becomes a cylinder S¹ × ℝ, and S³
becomes S³ × ℝ\_Q (Fig. 2, right). This axis cannot be read from the
tangential directions (which are readable as space). By the convention
of writing unreadable directions with imaginary units {[}S27{]}, (iQ)²
＝ −Q².

This is the origin of the Q axis of the author's spacetime equation

x² ＋ y² ＋ z² ＋ (it)² ＋ (iR)² ＋ (iQ)² ＝ 0.

The six terms split into 3 tangential (readable, space) ＋ 3 normal
(unreadable: time t, curvature R, amplitude direction Q). The
radial-projection framework ℝ⁶ → S⁵ {[}S42{]} and the internal-side
quaternion b ＝ R ＋ Q₁i ＋ Q₂j ＋ Q₃k of §11 of the previous paper
{[}S30{]} had placed this Q axis first. What this paper adds is the
origin: Q is the direction of the amplitude. With three Q's the normal
directions are five, t, R, Q₁, Q₂, Q₃, eight components in all:
octonions.

\begin{figure}
\centering
\pandocbounded{\includegraphics[keepaspectratio,alt={Where the amplitude lives. Left: on a circle in the plane, an amplitude placed radially is the same number as the radius of curvature R and cannot be distinguished from it. Right: on the cylinder S¹ × ℝ\_Q the amplitude lives in the unreadable direction Q, orthogonal both to the circle and to R.}]{figs/fig2_amplitude_radius_vs_Q_axis_en.png}}
\caption{Where the amplitude lives. Left: on a circle in the plane, an
amplitude placed radially is the same number as the radius of curvature
R and cannot be distinguished from it. Right: on the cylinder S¹ × ℝ\_Q
the amplitude lives in the unreadable direction Q, orthogonal both to
the circle and to R.}
\end{figure}

\textbf{Attribution.} That amplitude ≠ radius of curvature, and that the
amplitude lives on an orthogonal axis, become visible once a single
unreadable axis is added to the complex numbers. The four-dimensional
shell only made this easier to see. 【\(\mathbb{C}\) × ℝ】

\begin{center}\rule{0.5\linewidth}{0.5pt}\end{center}

\subsection{4. Reading the amplitude as imaginary: the stage of complex
numbers}\label{reading-the-amplitude-as-imaginary-the-stage-of-complex-numbers}

\subsubsection{4.1 Reading in the
exponent}\label{reading-in-the-exponent}

Read the relation u ＝ q₁⁻¹q₂ in its exponent. With w ＝ log u ＝ φ ＋
iθ, φ is the log-amplitude (amplitude ratio ρ ＝ e\^{}φ) and θ is the
phase. Evolution is a chain of multiplications, so in the exponent it
becomes addition and is linear. If amplitude and phase are written
separately as ψ ＝ A e\^{}\{iθ\}, the amplitude is treated by
multiplication and the phase inside the exponent, a different treatment,
and nonlinearity enters there. Madelung's hydrodynamic form {[}X11{]}
and Bohm's quantum potential ∇²A/A {[}X12{]} diverge at the nodes of the
wave function because amplitude and phase are separated midway and one
divides by the amplitude; as long as the evolution is carried out in the
exponent alone, e\^{}w never vanishes for finite w. Nodes arise only
when the sum is taken at the moment of readout. Conserved quantities as
phase-space area and Wick rotation are treated in {[}S16{]}, and linear
simplexification by square-quantity readout in {[}S26{]}.

\textbf{Attribution.} Linearity in the exponent, and divergence when a
logarithm is taken midway, are properties of complex numbers.
【\(\mathbb{C}\)】

\subsubsection{4.2 The duality of amplitude and phase: multiplication by
i exchanges
them}\label{the-duality-of-amplitude-and-phase-multiplication-by-i-exchanges-them}

In the exponent plane, multiplying w ＝ a ＋ iθ by i gives i·w ＝ −θ ＋
ia, and the amplitude exponent a and the phase θ are exchanged (Fig. 1).
Amplitude and phase are conjugate, and the duality is a 90° rotation of
the exponent space.

This one rotation appears in many places in physics under different
names. The conjugacy of number (amplitude²) and phase {[}X19{]}
{[}X20{]} {[}X21{]}. Time and frequency, position and wavenumber. The
Wick rotation connecting real decay e\^{}\{−βH\} and imaginary
oscillation e\^{}\{−iHt\} {[}X17{]} {[}X18{]}, and the Laplace variable
s ＝ σ ＋ iω of electric circuits. In the author's system, the
symplectic symmetry and Wick rotation obtained when conserved quantities
are read as phase-space area {[}S16{]} correspond to this.

Which axis is read as ``the imaginary part'' is the choice of where to
stop the axis in this rotation, and is not fixed on the side of the
components {[}S3{]}. Even the integer winding number of the phase is a
product of readout that arises only when an axis is chosen and divided
by 2π (wound onto a circle).

\begin{figure}
\centering
\pandocbounded{\includegraphics[keepaspectratio,alt={The exponent plane w ＝ a ＋ iθ. Multiplication by i gives i·w ＝ −θ ＋ ia and exchanges the amplitude exponent and the phase. Number--phase conjugacy, time and frequency, Wick rotation and the Laplace variable s ＝ σ ＋ iω are appearances of the same rotation.}]{figs/fig1_exponent_plane_duality_en.png}}
\caption{The exponent plane w ＝ a ＋ iθ. Multiplication by i gives i·w
＝ −θ ＋ ia and exchanges the amplitude exponent and the phase.
Number--phase conjugacy, time and frequency, Wick rotation and the
Laplace variable s ＝ σ ＋ iω are appearances of the same rotation.}
\end{figure}

\textbf{Attribution.} The duality itself is a property of complex
numbers. 【\(\mathbb{C}\)】

\subsubsection{4.3 Consequences of the imaginary reading
(identities)}\label{consequences-of-the-imaginary-reading-identities}

To read the amplitude exponent as imaginary is to place the amplitude
inside the same exponent as the phase, with an imaginary unit attached.
The consequences come in two steps.

\textbf{Indefinite metric.} If the amplitude A is read as iA, then (iA)²
＝ −A² and the squared norm can be negative. It was because only
positive real amplitudes were admitted that the norm was positive
definite and readable as a probability; once imaginary amplitudes are
admitted, the norm becomes an indefinite metric that is positive,
negative or zero. This is the same structure as the light cone, and the
nodes ψ ＝ 0 correspond to null points.

\textbf{Two sections of the complex angle.} For w ＝ φ ＋ iθ the
following hold (\texttt{verify\_\allowbreak{}identities.py} C0--C2).

\textbar cosh(w/2)\textbar² ＝ (cosh φ ＋ cos
θ)/2,　\textbar sinh(w/2)\textbar² ＝ (cosh φ − cos θ)/2

\textbar cosh(w/2)\textbar² ＋ \textbar sinh(w/2)\textbar² ＝ cosh
φ,　\textbar cosh(w/2)\textbar² − \textbar sinh(w/2)\textbar² ＝ cos θ

Written with the complex angle ω ＝ θ − iφ (u ＝ e\^{}\{iω\}), cosh(w/2)
＝ cos(ω/2) and sinh(w/2) ＝ i sin(ω/2), so the same formulas can be
written with cos and sin. Three things can be read from this (Fig. 3).

\begin{enumerate}
\def\labelenumi{\arabic{enumi}.}
\tightlist
\item
  \textbf{The section φ ＝ 0 is cos²(θ/2) and sin²(θ/2), with sum 1.}
  This is the form of the Born rule. That the projection weights of a
  discrete system with U² ＝ I are cos² and sin² (the same geometry as
  Malus's law), and that repeating a rotation by the angle θ ＝ πm/n
  closes the orbit with period n so that the fraction of one period
  spent on either projection side is the integer ratio m/n, have already
  been obtained in the author's system by {[}S19{]} {[}S20{]} {[}S21{]}
  from the equidistribution of the observer--system beat and the
  reproducing kernel. This paper rereads them as ``the φ ＝ 0 section of
  the complex angle.''
\item
  \textbf{The section θ ＝ 0 is cosh²(φ/2) and sinh²(φ/2), with sum cosh
  φ.} This is the Lorentz factor γ. That a boost is a rotation by an
  imaginary angle and that rapidities add has been the reading since
  Minkowski {[}X4{]}, and unifying rotations and boosts by a complex
  angle is standard in the spinor calculus of SL(2, \(\mathbb{C}\))
  {[}X5{]}.
\item
  \textbf{Nodes appear only on the section φ ＝ 0.} Even at θ ＝ π,
  \textbar cosh(w/2)\textbar² ＝ sinh²(φ/2), which does not vanish for φ
  ≠ 0. A node of zero amplitude appears only when one reads the single
  section φ ＝ 0.
\end{enumerate}

\begin{figure}
\centering
\pandocbounded{\includegraphics[keepaspectratio,alt={The two sections \textbar cosh(w/2)\textbar² and \textbar sinh(w/2)\textbar² of w ＝ log u ＝ φ ＋ iθ. The section φ ＝ 0 is cos²(θ/2) (the form of the Born rule). The sum is cosh φ ＝ γ independently of θ, and equals 1 only at φ ＝ 0.}]{figs/fig3_complex_angle_sections_gamma_en.png}}
\caption{The two sections \textbar cosh(w/2)\textbar² and
\textbar sinh(w/2)\textbar² of w ＝ log u ＝ φ ＋ iθ. The section φ ＝ 0
is cos²(θ/2) (the form of the Born rule). The sum is cosh φ ＝ γ
independently of θ, and equals 1 only at φ ＝ 0.}
\end{figure}

Since cosh x ＝ cos(ix), the circular and hyperbolic functions are the
same function read after a rotation by i. The cos² of quantum mechanics
and the γ of relativity are the real and the imaginary section of one
complex exponential; in terms of groups, SL(2, \(\mathbb{C}\)) is the
complexification of SU(2). It is not strange that either appears; it is
natural that both appear.

\textbf{Attribution.} The above are identities, and all lie at the stage
of complex numbers. 【\(\mathbb{C}\)】

\subsubsection{4.4 The Born-rule normalization and the denominator of
the previous
paper}\label{the-born-rule-normalization-and-the-denominator-of-the-previous-paper}

To make a general relation with φ ≠ 0 look like a probability, one must
divide by the sum cosh φ. \textbf{The one dimension discarded by that
division is the rapidity φ.} After the division the difference is cos θ
/ cosh φ, and the swing of the interference shrinks to 1/cosh φ ＝ √(1 −
v²). The moment one normalizes, the cause of that dilution becomes
unreadable.

The denominator of the previous paper's coupling coefficient k ＝ (1 −
u)/(1 ＋ u) is, with u ＝ e\^{}w (\texttt{verify\_\allowbreak{}identities.py}
C3--C5),

\textbar1 ＋ u\textbar² ＝ 1 ＋ 2e\^{}φ cos θ ＋ e\^{}\{2φ\} ＝
2e\^{}φ(cosh φ ＋ cos θ) ＝ 4e\^{}φ \textbar cosh(w/2)\textbar²,

with k ＝ −tanh(w/2) and \textbar k\textbar² ＝ (cosh φ − cos θ)/(cosh φ
＋ cos θ). \textbf{The denominator declared in §2.2 of the previous
paper as ``the composite q₁ ＋ q₂'' is the same division as the
Born-rule normalization.} Here appears the reason why the previous paper
demanded that division be done once, at the end, with the denominator
declared. If one divides midway, one cannot notice that the information
φ has disappeared.

And \textbar k\textbar{} ＝ 1 holds when cos θ ＝ 0, that is, at θ ＝
±π/2, \textbf{independently of the amplitude mismatch φ} (Fig. 4, left).
The previous papers {[}S27{]} {[}S30{]} identified the relation
\textbar k\textbar{} ＝ 1 with a relation on the light cone. This is
also the point where the map of coupling k and the map of distance 1/k
exchange and coincide (C10). The light cone is the self-dual point of
coupling and distance.

\begin{figure}
\centering
\pandocbounded{\includegraphics[keepaspectratio,alt={Left: \textbar k\textbar{} ＝ \textbar tanh(w/2)\textbar{} (log₁₀). The self-dual point \textbar k\textbar{} ＝ 1 lies at θ ＝ ±π/2, independently of φ. Right: the denominator \textbar1 ＋ u\textbar² ＝ 2e\^{}φ(cosh φ ＋ cos θ) ＝ 4e\^{}φ\textbar cosh(w/2)\textbar².}]{figs/fig4_k_selfdual_denominator_en.png}}
\caption{Left: \textbar k\textbar{} ＝ \textbar tanh(w/2)\textbar{}
(log₁₀). The self-dual point \textbar k\textbar{} ＝ 1 lies at θ ＝
±π/2, independently of φ. Right: the denominator \textbar1 ＋ u\textbar²
＝ 2e\^{}φ(cosh φ ＋ cos θ) ＝ 4e\^{}φ\textbar cosh(w/2)\textbar².}
\end{figure}

\textbf{Attribution.} The identity for the denominator and the self-dual
point lie at the stage of complex numbers. 【\(\mathbb{C}\)】

\subsubsection{4.5 In complex numbers, amplitude and phase cannot be
distinguished}\label{in-complex-numbers-amplitude-and-phase-cannot-be-distinguished}

Now apply the decision rule. If in complex numbers the amplitude
exponent is read as imaginary, one gets e\^{}\{i(a ＋ θ)\}: a is added
to θ and the modulus stays 1 (C9). An amplitude read as imaginary cannot
be told from a shift of the phase by π/2. There is only one imaginary
unit. What remains of the imaginary reading at the stage of complex
numbers is only the sign reversal of the square (the sign of the Wick
rotation).

\textbf{Attribution.} That ``amplitude and phase cannot be
distinguished'' is itself a property of the stage of complex numbers.
The distinction first arises at the next stage. 【\(\mathbb{C}\)】

\subsubsection{4.6 Which section to read, and where to
divide}\label{which-section-to-read-and-where-to-divide}

Of the two sections of §4.3, quantum mechanics chooses φ ＝ 0 and calls
it ``probability,'' and relativity looks at θ ＝ 0 and calls it ``the
upper limit of velocity.'' Up to here it is a result of algebra, not
physics. What enters as physics is only the choice: which section to
read, and where to divide in order to normalize.

If one tries to derive this choice from dynamics, the dynamics itself is
already written on some section, so one would be explaining the choice
by presupposing it. As long as the observer is a part of the system and
the readout is included in what is read out, this self-reference cannot
be avoided. Therefore \textbf{which axis is read as the imaginary part,
and where one divides, are placed as axioms.} Existing physics places
the Born rule and the invariance of the speed of light as two separate
axioms, but in the organization of this section both come from a single
choice, namely which section of the same complex exponential to read and
where to divide. It cannot be derived, but the two axioms are unified
into one.

The view that the choice of imaginary unit selects the kind of mechanics
is closest to Kisil {[}X16{]}, who divided the representations of the
Heisenberg group into elliptic (i² ＝ −1, quantum mechanics), parabolic
(ε² ＝ 0, classical mechanics) and hyperbolic (j² ＝ ＋1, hyperbolic
mechanics). A quantum mechanics that reads the amplitude in split form
was proposed by Khrennikov {[}X13{]} {[}X14{]}, with the criticism that
negative probabilities appear {[}X15{]}. The standpoint of reading
probabilities as a counting of rational numbers m/n is found in Palmer's
invariant set theory {[}X22{]} {[}X23{]} and 't Hooft's cellular
automaton interpretation {[}X24{]}; in the author's system, the finite
order U\^{}n ＝ I {[}S2{]} {[}S22{]} and rational locking {[}S23{]}
{[}S24{]} correspond to it.

\textbf{Attribution.} That the choice is an axiom is common to all
stages. 【\(\mathbb{C}\)】

\begin{center}\rule{0.5\linewidth}{0.5pt}\end{center}

\subsection{5. The stage of quaternions: amplitude and phase are
distinguished, and space is
born}\label{the-stage-of-quaternions-amplitude-and-phase-are-distinguished-and-space-is-born}

\subsubsection{5.1 Giving a separate unit}\label{giving-a-separate-unit}

In quaternions the amplitude can be given an imaginary unit m (m ⊥ n)
different from that of the phase. Since the exponent of q ＝ e\^{}\{ma
＋ nθ\} is pure imaginary, \textbar q\textbar{} ＝ 1 stays, and
everything is a rotation on S³. The amplitude moves, but the norm does
not. ``The overall scale cannot be read'' {[}S3{]} becomes a structure
rather than an axiom. Since m and n are different, orthogonal
directions, amplitude and phase can be distinguished.

\textbf{Attribution.} In complex numbers there was only one unit and no
distinction (§4.5). That they can be distinguished is what quaternions
add. 【\(\mathbb{H}\)】

\subsubsection{5.2 Non-commutativity generates a third
direction}\label{non-commutativity-generates-a-third-direction}

The rotations by m and by n do not commute. For unit pure imaginary m ⊥
n (\texttt{verify\_\allowbreak{}identities.py} C7),

e\^{}\{ma\} e\^{}\{nθ\} − e\^{}\{nθ\} e\^{}\{ma\} ＝ 2 sin a sin θ (m ×
n).

The two imaginary readings, of amplitude and of phase, generate through
the order of multiplication a third direction m × n, and the three fill
the three-dimensional imaginary part, that is, space (Fig. 5). The axis
of the composite rotation reverses the sign of its m × n component with
the order of multiplication. This difference of order is the left and
right division of the previous paper {[}S27{]}, the choice of the
observer's frame. Moreover e\^{}\{ma ＋ nθ\} and e\^{}\{ma\} e\^{}\{nθ\}
are different things (the Baker--Campbell--Hausdorff correction, C8), so
which of them is called ``the state combining amplitude and phase'' also
requires a declaration of order (§2.2).

Viewed in quaternions, the duality of amplitude and phase (§4.2), the
rotation that exchanges the amplitude axis m and the phase axis n, is a
90° spatial rotation about the third axis m × n.~The duality that was
``a rotation of the exponent plane'' at the stage of complex numbers
becomes, in quaternions, the same thing as one rotation of space.

\begin{figure}
\centering
\pandocbounded{\includegraphics[keepaspectratio,alt={Quaternions: the rotations about the amplitude axis m and the phase axis n do not commute, and the commutator e\{ma\}e\{nθ\} − e\{nθ\}e\{ma\} ＝ 2 sin a sin θ (m × n) generates a third direction. The axis of the composite rotation reverses the sign of its m × n component with the order.}]{figs/fig5_quaternion_commutator_third_axis_en.png}}
\caption{Quaternions: the rotations about the amplitude axis m and the
phase axis n do not commute, and the commutator
e\textsuperscript{\{ma\}e}\{nθ\} − e\textsuperscript{\{nθ\}e}\{ma\} ＝ 2
sin a sin θ (m × n) generates a third direction. The axis of the
composite rotation reverses the sign of its m × n component with the
order.}
\end{figure}

\textbf{Attribution.} The third direction comes from non-commutativity.
Complex numbers do not have it. 【\(\mathbb{H}\)】

\subsubsection{5.3 The k of the previous paper is the modified Rodrigues
parameter}\label{the-k-of-the-previous-paper-is-the-modified-rodrigues-parameter}

The coupling coefficient of the previous paper is, for a unit quaternion
u ＝ e\^{}\{nα\}, k ＝ (1 ＋ u)⁻¹(1 − u) ＝ −n tan(α/2) ＝ −Im u / (1 ＋
Re u) (C6). This is the \textbf{modified Rodrigues parameter} (MRP) with
its sign reversed. The MRP is the stereographic coordinate of a unit
quaternion, long used in attitude dynamics in aerospace {[}X29{]}
{[}X30{]}, and introduced to computer graphics and vision by Terzakis et
al.~{[}X28{]}. What the previous paper called ``the central projection
of the square root of the relative phase'' is this known map.

The existing MRP is used to represent a rotational attitude within
three-dimensional space. The previous paper used the same map to read
the relation between states as a position, without placing a background
space. The rediscovery of the map and its positioning as a readout
principle are written separately. The classical maps from two complex
numbers to a three-dimensional vector (the Kustaanheimo--Stiefel
transformation {[}X31{]}, the Hopf map {[}X32{]}) contain the choice of
a special direction in their definition and differ in respect of
anonymity.

The root of k being a three-dimensional vector lies in the third
direction of §5.2. The core of the previous paper {[}S27{]}, that two
standing waves can intrinsically memorize a three-dimensional position,
can be read as a consequence of the non-commutativity of the two
imaginary readings, of amplitude and of phase.

\textbf{Attribution.} k as a three-dimensional vector lies at the stage
of quaternions. 【\(\mathbb{H}\)】

\subsubsection{5.4 (1, 3), and the light cone with an
orientation}\label{and-the-light-cone-with-an-orientation}

The previous paper {[}S27{]} showed that the scalar part q₀² −
\textbar v\textbar² of the true square, taken without the conjugate, has
signature (1, 3), and that θ ＝ 45° is the light cone {[}S28{]}. In the
light of §4.3, the sign reversal (the light cone of 1 ＋ 1) is already
at the stage of complex numbers, and what quaternions add is that the
light cone is not a pair of points but \textbf{a cone with an
orientation n}. The attribution ``the light cone comes from the true
square,'' written in the previous paper, is refined in this way. That
\textbar k\textbar{} ＝ 1 is a relation on the light cone is itself at
the stage of complex numbers, as in §4.4; quaternions give it an
orientation.

\textbf{Attribution.} The sign is 【\(\mathbb{C}\)】; the oriented cone
is 【\(\mathbb{H}\)】.

\subsubsection{5.5 Why condensates cannot be made with complex numbers
and can be made with
quaternions}\label{why-condensates-cannot-be-made-with-complex-numbers-and-can-be-made-with-quaternions}

The author was stuck for a long time on the problem of making a local
binding structure (a condensate) from an anonymous set of complex
numbers without supplying a background or connectivity. The complex
relative phase u ＝ e\^{}\{i(θ\_b − θ\_a)\} is a single angle, and the
product is commutative. It is not determined in which direction several
relations spread and how they combine to make a three-dimensional local
structure. In quaternions the relative phase u ＝ e\^{}\{nα\} has a
direction n, and for the directions of two relations there is the
non-commutative structure ab − ba ＝ 2a × b, so the distinction of
directions, parallelism, and the construction of triangles and
tetrahedra can be considered. The convex hull and the tetrahedral
crystallization of the previous paper {[}S30{]} come from here.

That three quaternions q₁, q₂, q₃ do not fit into a single complex phase
plane can be tested with the normalization pᵢ ＝ qᵢ/\textbar qᵢ\textbar,
gᵢⱼ ＝ Re(p̄ᵢpⱼ) and the discriminant Δ₃ ＝ det{[}gᵢⱼ{]} ＝ ¼\textbar uv
− vu\textbar² (u ＝ q₁⁻¹q₂, v ＝ q₁⁻¹q₃). If Δ₃ ＝ 0 the directions of
the relative phases are collinear and the configuration is complex-like;
if Δ₃ \textgreater{} 0 there are two independent directions and it is
quaternionic. The definition of a condensate (a common curvature scale,
locking) is treated in the next paper.

A distinction is made here. Reading out localized condensates, a centre
of mass and a mass-like measure from a vacuum of complex standing waves
has been done in {[}S31{]} {[}S32{]}. What is said here to be impossible
is not readout but \textbf{autonomous binding} (which state binds to
which being decided without being given from outside). That the birth
conditions of space, matter and the clock do not coincide {[}S33{]} is
related to this distinction.

\textbf{Attribution.} Relations with a direction, and the possibility of
condensates, lie at the stage of quaternions. 【\(\mathbb{H}\)】

\begin{center}\rule{0.5\linewidth}{0.5pt}\end{center}

\subsection{6. The stage of octonions: mass-type and charge-type
amplitudes are
distinguished}\label{the-stage-of-octonions-mass-type-and-charge-type-amplitudes-are-distinguished}

\subsubsection{6.1 Two elements are associative; the associator stands
with
three}\label{two-elements-are-associative-the-associator-stands-with-three}

The subalgebra generated by any two octonions is associative (Artin's
theorem {[}X36{]}). However much two elements are multiplied, one cannot
get out of the quaternions. The associator {[}a, b, c{]} ＝ (ab)c −
a(bc) can first be tested with three elements, and when the true-square
closure is also imposed, the smallest set that closes as octonions is
four (five if the linear closure is imposed as well) {[}S30{]} {[}S4{]}.
Together with the fact that number systems with division are limited to
1, 2, 4 and 8 dimensions (Hurwitz {[}X37{]}, Baez {[}X27{]}), the
octonions are the stage that can first be described only when the order
of products of triples is included.

\subsubsection{6.2 The type of amplitude splits in
two}\label{the-type-of-amplitude-splits-in-two}

When the imaginary units of amplitude are two (m in the direction of R
and l in the direction of Q) and, together with the phase unit n, form
three independent imaginary directions, the associator first becomes
non-zero here. In quaternions, even if R and Q are given separate units,
two imaginary units always fall within an associative subalgebra, and
which one is called mass is only a renaming. In octonions the triple R,
Q, phase becomes non-associative, and \textbf{the R-type amplitude and
the Q-type amplitude become algebraically different things.}

This gives the question ``why can mass and charge be measured as
distinct; are they separate entities, or different readouts of the same
thing?'' the algebraic answer ``because the associator first stands with
the third imaginary reading.'' The very fact that they are
distinguishable is the non-associativity of three imaginary directions.

\textbf{Attribution.} The distinction of two kinds of amplitude lies at
the stage of octonions. 【\(\mathbb{O}\)】

\subsubsection{6.3 The split form and \textbar a\textbar{} ＝
\textbar b\textbar{}}\label{the-split-form-and-a-b}

Reading all the internal-side amplitudes (R, Q₁, Q₂, Q₃) as imaginary is
exactly choosing the split octonions of §11 of the previous paper
{[}S30{]} (γ ＝ ＋1, signature (4, 4)). With the Cayley--Dickson
construction (a, b), a the spacetime-side and b the internal-side
quaternion, the null condition N(o) ＝ \textbar a\textbar² −
\textbar b\textbar² ＝ 0 is \textbar a\textbar{} ＝ \textbar b\textbar,
that is,

x² ＋ y² ＋ z² ＋ t² ＝ R² ＋ Q₁² ＋ Q₂² ＋ Q₃².

Since the norm of the ordinary octonions is positive definite and has no
non-trivial null elements, the split form is required for the closure to
be non-trivial. This is the same logic as Σxₙ² ＝ 0 having only the
trivial solution over the reals and requiring complex numbers {[}S4{]}
{[}S28{]}.

\subsubsection{6.4 Reading B: mass ＝ spacetime curvature R, charge ＝
internal curvature
Q}\label{reading-b-mass-spacetime-curvature-r-charge-internal-curvature-q}

In §3.4 the amplitude was placed on the Q axis. Here two readings stand
side by side.

\begin{itemize}
\tightlist
\item
  \textbf{Reading A.} The amplitude lives in Q; mass ＝ amplitude,
  charge ＝ its direction.
\item
  \textbf{Reading B.} Mass ＝ the spacetime curvature R itself; charge
  ＝ the internal curvature Q itself. Charge does not affect the
  spacetime curvature and affects only the internal curvature.
\end{itemize}

This paper adopts B, for three reasons. First, that mass curves
spacetime is general relativity itself, and ``the source of gravity is
the non-uniformity of the norm'' of §6.3 of the previous paper {[}S30{]}
holds as it is when reread as ``the non-uniformity of R.'' Second,
charge as internal geometry is the Kaluza--Klein type of reading. Third,
it agrees with the author's internal records: the kind of a particle is
read by the sign vector of tRQ₁Q₂Q₃, and the charge comes from the signs
of Q {[}S34{]} {[}S35{]} {[}S36{]}; the structure of generations
{[}S37{]}; the hierarchy of masses appears as a phase of R; and changing
the normalization reference from mass to charge does not change the
generation of states {[}S25{]}.

The dictionary of Reading B repeats the same 1 ＋ 3 on the spacetime
side and on the internal side (Table 1).

\textbf{Table 1} Dictionary of the spacetime side and the internal side
(Reading B).

{\def\LTcaptype{none} % do not increment counter
\begin{longtable}[]{@{}
  >{\raggedright\arraybackslash}p{(\linewidth - 6\tabcolsep) * \real{0.2500}}
  >{\raggedright\arraybackslash}p{(\linewidth - 6\tabcolsep) * \real{0.2500}}
  >{\raggedright\arraybackslash}p{(\linewidth - 6\tabcolsep) * \real{0.2500}}
  >{\raggedright\arraybackslash}p{(\linewidth - 6\tabcolsep) * \real{0.2500}}@{}}
\toprule\noalign{}
\begin{minipage}[b]{\linewidth}\raggedright
Side
\end{minipage} & \begin{minipage}[b]{\linewidth}\raggedright
Quaternion
\end{minipage} & \begin{minipage}[b]{\linewidth}\raggedright
Scalar part
\end{minipage} & \begin{minipage}[b]{\linewidth}\raggedright
Vector part
\end{minipage} \\
\midrule\noalign{}
\endhead
\bottomrule\noalign{}
\endlastfoot
Spacetime side & a ＝ t ＋ xi ＋ yj ＋ zk & t ＝ time & (x, y, z) ＝
space \\
Internal side & b ＝ R ＋ Q₁i ＋ Q₂j ＋ Q₃k & R ＝ mass (spacetime
curvature) & (Q₁, Q₂, Q₃) ＝ charge (internal curvature) \\
\end{longtable}
}

The amplitude is the norm of the internal side, \textbar b\textbar{} ＝
√(R² ＋ ΣQ²), and the null condition \textbar a\textbar{} ＝
\textbar b\textbar{} is ``the inventory of the spacetime side ＝ the
inventory of the internal side.'' Hypothesis 1 of the previous paper,
``mass ∝ \textbar q\textbar²,'' is refined in Reading B to ``mass ∝ the
scalar part R of the internal side,'' and the Q part of the amplitude is
read not as mass but as charge. This paper states this explicitly as a
correction of §6 of the previous paper.

Two readings follow from Reading B (not yet checked). Charge is the
\textbf{direction and sign} of the internal 3-vector; the sign vector is
discrete (16 for one generation {[}S34{]} {[}S36{]}), so it is
quantized. Mass is the \textbf{magnitude} of a scalar and can take
continuous values. This agrees with charge appearing only in integer
multiples of e while mass has a hierarchy but no visible rule of
quantization. Also, the elementary charge read as a transmittance, 1 − R
＝ sin²θ ＝ √(4πα) ＝ e\_HL {[}S17{]} {[}S18{]} {[}S38{]} {[}S39{]}
{[}S40{]} {[}S41{]}, says that the \textbf{magnitude} of charge is
universal (the length of the internal vector is fixed), and this is read
separately from the sign and direction being the choice of Q.

The k of the previous paper is built from the spacetime side a, so its
vector part is the direction of space, not charge. To read charge one
needs the relative phase u\_b ＝ b₁⁻¹b₂ of the internal side b and
k\_internal ＝ (1 ＋ u\_b)⁻¹(1 − u\_b). How its scalar part (the ratio
of R ＝ the mass ratio) and its vector part (the difference of the
directions of Q ＝ the relation of charges) are connected with the
spacetime-side k and with \textbar a\textbar{} ＝ \textbar b\textbar{}
is what decides between Readings A and B, and is the next task (§9).

\textbf{Attribution.} Which normal direction receives the name mass and
which the name charge is a choice of readout; the components carry no
colour. What decides is whether a general-relativity-type readout on the
R side and a gauge-type readout on the Q side come out without
contradiction. That this distinction holds algebraically is at the stage
of octonions. 【\(\mathbb{O}\)】

\subsubsection{6.5 A stage without
precedent}\label{a-stage-without-precedent}

By Dyson's threefold way {[}X41{]}, the ensembles of matrices (the
symmetry classes of random matrices) are only three, ℝ, \(\mathbb{C}\)
and \(\mathbb{H}\); there is none for the octonions. Up to quaternions
there are precedents of exactly solvable many-body systems, the log gas,
Dyson's Brownian motion and the Sutherland model, but there is no
precedent for the many-body dynamics of octonions. The correspondence
between octonions and the Standard Model has a lineage since Gürsey--Tze
{[}X38{]} {[}X39{]} {[}X40{]}; this paper stands in the position of
adding to it the reading of two types of amplitude. Dynamics is outside
the scope of this paper and is treated in the next paper (§9).

\begin{center}\rule{0.5\linewidth}{0.5pt}\end{center}

\subsection{7. The attribution ledger}\label{the-attribution-ledger}

Sorting the claims by the decision rule gives Table 2 (Fig. 6).

\textbf{Table 2} The attribution ledger.

{\def\LTcaptype{none} % do not increment counter
\begin{longtable}[]{@{}
  >{\raggedright\arraybackslash}p{(\linewidth - 4\tabcolsep) * \real{0.3333}}
  >{\raggedright\arraybackslash}p{(\linewidth - 4\tabcolsep) * \real{0.3333}}
  >{\raggedright\arraybackslash}p{(\linewidth - 4\tabcolsep) * \real{0.3333}}@{}}
\toprule\noalign{}
\begin{minipage}[b]{\linewidth}\raggedright
Already in \(\mathbb{C}\)
\end{minipage} & \begin{minipage}[b]{\linewidth}\raggedright
Added by \(\mathbb{H}\)
\end{minipage} & \begin{minipage}[b]{\linewidth}\raggedright
Added by \(\mathbb{O}\)
\end{minipage} \\
\midrule\noalign{}
\endhead
\bottomrule\noalign{}
\endlastfoot
Duality of amplitude and phase (90° rotation of the exponent plane) §4.2
& Amplitude and phase distinguishable (m ⊥ n) §5.1 & Non-associativity:
the associator stands with 3 elements, and with closure imposed, 4
§6.1 \\
Imaginary amplitude → indefinite metric, light cone §4.3 &
Non-commutativity {[}m, n{]} ＝ 2m × n → third direction ＝ space §5.2 &
R-type and Q-type amplitudes algebraically distinct → distinction of
mass and charge §6.2 \\
Two sections of the complex angle: cos² (Born rule) and cosh φ ＝ γ §4.3
& Sign (1, 3) of the true square, light cone with an orientation §5.4 &
Split form (4, 4) and \textbar a\textbar{} ＝ \textbar b\textbar: the
spacetime equation §6.3 \\
Self-dual point \textbar k\textbar{} ＝ 1 (θ ＝ ±π/2, independent of φ)
§4.4 & Two standing waves memorize a position §5.3 & The order of
sequential binding changes the result (history) §6.1 \\
Σz ＝ 0 and Σz² ＝ 0, U\^{}n ＝ I, counting of m/n §4.6 & Condensates
can be made (complex numbers have no direction) §5.5 & No Dyson ensemble
exists (no precedent) §6.5 \\
The division trap, linearity in the exponent §4.1, §4.4 & Three-body
discriminant Δ₃ ＝ ¼\textbar uv − vu\textbar² §5.5 & \\
Amplitude ≠ radius of curvature (visible on the cylinder S¹ × ℝ) §3 &
& \\
\end{longtable}
}

\begin{figure}
\centering
\pandocbounded{\includegraphics[keepaspectratio,alt={The attribution ledger: already in \textbackslash mathbb\{C\} / added by \textbackslash mathbb\{H\} / added by \textbackslash mathbb\{O\}.}]{figs/fig6_attribution_ledger_en.png}}
\caption{The attribution ledger: already in \(\mathbb{C}\) / added by
\(\mathbb{H}\) / added by \(\mathbb{O}\).}
\end{figure}

Three examples in which the author nearly misattributed during the check
are recorded. First, the orthogonality of amplitude and curvature was
about to be written as a property of the four-dimensional shell, but it
is visible on the cylinder S¹ × ℝ. Second, one is tempted to read γ as a
consequence of the quaternionic signature (1, 3), but it lies in the
complex angle; what quaternions add is not the sign but the three
directions of space. Third, the previous paper {[}S27{]} wrote that the
light cone ``comes from the true square,'' but the relation
\textbar k\textbar{} ＝ 1 is in the complex numbers, and what
quaternions added is the orientation of the cone.

\begin{center}\rule{0.5\linewidth}{0.5pt}\end{center}

\subsection{8. Comparison with prior
work}\label{comparison-with-prior-work}

Each of the parts of this paper has prior work that gave it a physical
interpretation. The position of this paper is to have put them together
into one line, and in the right-hand columns of the ledger.

{\def\LTcaptype{none} % do not increment counter
\begin{longtable}[]{@{}
  >{\raggedright\arraybackslash}p{(\linewidth - 4\tabcolsep) * \real{0.3333}}
  >{\raggedright\arraybackslash}p{(\linewidth - 4\tabcolsep) * \real{0.3333}}
  >{\raggedright\arraybackslash}p{(\linewidth - 4\tabcolsep) * \real{0.3333}}@{}}
\toprule\noalign{}
\begin{minipage}[b]{\linewidth}\raggedright
Part
\end{minipage} & \begin{minipage}[b]{\linewidth}\raggedright
Prior work
\end{minipage} & \begin{minipage}[b]{\linewidth}\raggedright
Difference from this paper
\end{minipage} \\
\midrule\noalign{}
\endhead
\bottomrule\noalign{}
\endlastfoot
Phase ＝ rotation in a plane of space & Hestenes' real Dirac theory
{[}X6{]} {[}X7{]} {[}X8{]}: the i of the Dirac equation is the generator
of rotation in the spatial plane corresponding to the spin direction. β
(Yvon {[}X9{]}, Takabayasi {[}X10{]}) is the angle of the duality
rotation & Keeps the amplitude ρ real and separate from the phase. No
imaginary reading of the amplitude \\
Boost ＝ rotation by an imaginary angle, γ ＝ cosh θ & Minkowski
{[}X4{]}, Penrose--Rindler {[}X5{]} & No identification with the
Born-rule normalization \\
Which imaginary unit is read selects the mechanics & Kisil {[}X16{]}
(elliptic, parabolic, hyperbolic), Khrennikov {[}X13{]} {[}X14{]},
criticism {[}X15{]} & The unit is chosen per theory; not two sections of
one and the same object \\
Real exponent ↔ imaginary exponent ＝ temperature ↔ time & KMS,
Connes--Rovelli {[}X17{]}, imaginary time {[}X18{]} & Limited to the
exchange of statistics and quantum theory \\
Number--phase conjugacy & Dirac {[}X19{]}, Susskind--Glogower {[}X20{]},
Pegg--Barnett {[}X21{]} & Conjugacy of operators; no reading as a
rotation of the exponent space \\
cos² ＝ counting of m/n, Bell by discreteness & Palmer {[}X22{]}
{[}X23{]}, 't Hooft {[}X24{]} & From fractal invariant sets and
non-computability, not from the finite order U\^{}n ＝ I \\
The map k ＝ −n tan(α/2) & Modified Rodrigues parameters {[}X28{]}
{[}X29{]} {[}X30{]} & Representation of a rotational attitude, not a
readout of position \\
Reconstructing geometry from relations & The Penrose--Szabados line
{[}X33{]} {[}X34{]} {[}X35{]}: three-space and Minkowski spacetime
reconstructed from the symmetries and observables of quantum systems,
without a background physical space and without using an evolution
equation & Starts from known symmetry groups (SU(2), E(3), Poincaré) and
the observables momentum and angular momentum. Not a construction from
the relative relations of anonymous stationary states \\
Octonions and the Standard Model & Singh {[}X40{]} (matrix dynamics, the
exceptional Jordan algebra, 1/137), Furey {[}X38{]} {[}X39{]} (the
Standard Model groups from the agreement of the action of symmetry
transformations) & Not the minimal construction from stationary states
in which not even a progressing phase is placed. This paper adds the two
types of amplitude \\
Three-dimensional space does not emerge uniquely from the quantum
structure & Stoica {[}X1{]} & This paper takes the same standpoint (a
choice of reading enters). Definition and geometry of quaternions:
{[}X25{]} {[}X26{]} \\
\end{longtable}
}

Only this paper has the following four points. That cosh φ ＝ γ is the
one dimension discarded by the Born-rule normalization, and that the
Born rule and the invariance of the speed of light are two sections of
one axiom (§4.4, §4.6). That the self-dual point of the amplitude--phase
duality is the light cone (§4.4). That raising the imaginary reading
through \(\mathbb{C}\) → \(\mathbb{H}\) → \(\mathbb{O}\) produces first
the distinction from the phase and then the distinction of mass-type and
charge-type by the associator (§5, §6). That the amplitude is housed in
the normal direction Q of S³ (§3.4).

The literature on dynamics (the Kuramoto model {[}X42{]}, the Kuramoto
model with complexified phase {[}X43{]} {[}X44{]}, the author's
Kuramoto-type phase coupling {[}S43{]}, the linear instability of
inflation {[}S44{]} {[}S45{]} {[}S46{]} {[}S47{]}) is only cited in this
paper; its results are not used.

\begin{center}\rule{0.5\linewidth}{0.5pt}\end{center}

\subsection{9. Classification of claims, verification, and falsification
conditions}\label{classification-of-claims-verification-and-falsification-conditions}

\subsubsection{9.1 Classification of
claims}\label{classification-of-claims}

{\def\LTcaptype{none} % do not increment counter
\begin{longtable}[]{@{}
  >{\raggedright\arraybackslash}p{(\linewidth - 2\tabcolsep) * \real{0.5000}}
  >{\raggedright\arraybackslash}p{(\linewidth - 2\tabcolsep) * \real{0.5000}}@{}}
\toprule\noalign{}
\begin{minipage}[b]{\linewidth}\raggedright
Claim
\end{minipage} & \begin{minipage}[b]{\linewidth}\raggedright
Class
\end{minipage} \\
\midrule\noalign{}
\endhead
\bottomrule\noalign{}
\endlastfoot
The identities of §4.3 (two sections, sum ＝ cosh φ, difference ＝ cos
θ), the denominator and the formula for \textbar k\textbar² of §4.4, the
commutator of §5.2, the agreement with the MRP of §5.3, the merging of
§4.5, Artin's theorem of §6.1 & Mathematically identical (checked
numerically by \texttt{verify\_\allowbreak{}identities.py} C0--C10) \\
The amplitude cannot live radially and lives in the orthogonal imaginary
direction (§3); which section to read and where to divide is an axiom
(§4.6); Reading B (§6.4); the attribution ledger (§7) & Choice of
reading \\
Why charge is quantized and mass looks continuous (§6.4); that
k\_internal settles Reading B (§6.4) & Hypothesis \\
\end{longtable}
}

\subsubsection{9.2 What was checked}\label{what-was-checked}

\texttt{図表/\allowbreak{}verify\_\allowbreak{}identities.py} checks 11 identities on a fixed grid
(θ ∈ {[}−π, π{]} × φ ∈ {[}−3, 3{]}, 721 × 301 points), and all agree to
within 10⁻¹⁰ (\texttt{results/\allowbreak{}checks.txt}). No random numbers are used.
The quaternion identities (C6--C8) were computed directly from products
of unit quaternions.

\subsubsection{9.3 Falsification of the
attribution}\label{falsification-of-the-attribution}

The application of the decision rule itself has the form of a
falsification.

\begin{enumerate}
\def\labelenumi{\arabic{enumi}.}
\tightlist
\item
  If a readout of 3 directions can be built with complex numbers alone,
  ``the third direction is added by quaternions'' (§5.2) is wrong.
\item
  If the distinction of mass and charge comes out algebraically with
  quaternions alone, ``the distinction of two kinds of amplitude is
  added by octonions'' (§6.2) is wrong.
\item
  If a reading can be built that places the amplitude radially and keeps
  a common curvature scale inside the same condensate, §3.3 falls.
\item
  If, on computing the internal-side k\_internal, Reading A gives a
  general-relativity-type readout on the R side and a gauge-type readout
  on the Q side without contradiction, Reading B falls.
\end{enumerate}

\subsubsection{9.4 Next tasks}\label{next-tasks}

\begin{enumerate}
\def\labelenumi{\arabic{enumi}.}
\tightlist
\item
  Write the internal-side k\_internal ＝ (1 ＋ u\_b)⁻¹(1 − u\_b) and
  check how its scalar part and vector part are connected with the
  spacetime-side k and with \textbar a\textbar{} ＝ \textbar b\textbar{}
  (settling Reading B).
\item
  The definition of a condensate (Δ₃, a common curvature scale,
  dynamical locking).
\item
  The declaration of whether the state combining amplitude and phase in
  the complex angle is written as e\^{}\{ma ＋ nθ\} or as e\^{}\{ma\}
  e\^{}\{nθ\}; left and right division; the order of brackets in
  octonions.
\item
  Dynamics: the base family (the complexified Kuramoto model {[}X43{]}
  {[}X44{]}, the author's Kuramoto-type coupling {[}S43{]}); the form of
  the function of the relative phase (the first harmonic (u − u⁻¹)/2, or
  the all-odd-harmonic (1 − u)/(1 ＋ u)); the minimal system (octonions
  need 4, 5 if linear closure is imposed as well, 8 or more if isotropy
  is imposed); two-body reactions with the observer as a distant
  absorber {[}S48{]}, the exchange of harmonics of localized waves
  {[}S49{]}, and the discrimination of pairs that are close but do not
  lock {[}S50{]}.
\item
  The numerical experiments V1--V5 of §13 of the previous paper
  {[}S30{]}. They were not run in this paper.
\end{enumerate}

\begin{center}\rule{0.5\linewidth}{0.5pt}\end{center}

\subsection{10. Conclusion}\label{conclusion}

We began from the naive question of where the amplitude lives. On a
circle in the plane the amplitude and the radius of curvature are the
same number and cannot be distinguished, but the radius of curvature is
a quantity that can be read only by the area of a closed loop, and the
amplitude cannot live there. The amplitude lives in an unreadable
direction orthogonal both to the circle and to the radius of curvature,
and that is the origin of the Q axis of the author's spacetime equation.

Reading the amplitude in the unreadable direction as imaginary, there
appear already at the stage of complex numbers the duality of amplitude
and phase (the 90° rotation of the exponent plane), the indefinite
metric and the light cone, the two sections cos² (the form of the Born
rule) and cosh φ ＝ γ (the Lorentz factor) of the same complex
exponential, the self-dual point \textbar k\textbar{} ＝ 1 ＝ light cone
of coupling and distance, and the fact that the denominator \textbar1 ＋
u\textbar² of the previous paper is the same division as the Born-rule
normalization. Which section to read and where to divide is an axiom,
and the Born rule and the invariance of the speed of light are two
sections of that one axiom.

In complex numbers, an amplitude read as imaginary and the phase cannot
be distinguished. Quaternions distinguish them with a separate unit m ⊥
n, and the non-commutativity {[}m, n{]} ＝ 2m × n generates a third
direction, that is, space. The k of the previous paper is the modified
Rodrigues parameter, and the root of two standing waves being able to
memorize a position lies here. Octonions split the type of amplitude
into two, R and Q, and by the associator of triples make mass-type and
charge-type amplitudes algebraically different things. The same 1 ＋ 3
repeats on the spacetime side a ＝ t ＋ xyz and the internal side b ＝ R
＋ Q₁Q₂Q₃, and they are read as mass ＝ spacetime curvature R and charge
＝ internal curvature Q.

The attribution ledger puts what each stage adds into a form falsifiable
one stage down. The answer to the question whether we understood complex
numbers correctly is this: ``What lies at the stage of complex numbers
is more than we thought; what quaternions and octonions add is more
limited than we thought; and that limited part is space, and the
distinction of mass and charge.''

\begin{center}\rule{0.5\linewidth}{0.5pt}\end{center}

\subsection{References}\label{references}

\subsubsection{Self-citations (Concept DOI,
10.5281/zenodo.〔number〕)}\label{self-citations-concept-doi-10.5281zenodo.number}

{[}S1{]} N. Kihara, ``Anonymous Equal-Amplitude Composite-Wave Model:
Basic Axiom System v9,'' 2026-07-24. Concept DOI
10.5281/zenodo.21315735.

{[}S2{]} N. Kihara, ``Symmetry Generation from Zero Closure, Finite
Order, and Self-Consistent Geometry: The Single External Parameter N and
the Remaining Tasks of Generalization and Dynamics,'' 2026-08-20.
Concept DOI 10.5281/zenodo.22028072.

{[}S3{]} N. Kihara, ``A Catalogue of Symmetries and Magic Numbers for
Reading Out Spacetime from a State Space: The Answer-Check Standard of a
Research Programme Taking Anonymity and Self-Consistency as Axioms,''
2026-10-18. Concept DOI 10.5281/zenodo.23444451.

{[}S4{]} N. Kihara, ``An Axiomatic Framework for Physical Symmetry
Selection by Anonymous Partial Projection: Conditional Identification of
E8 in an Eight-Component Two-Quartet Branch,'' 2026-07-24. Concept DOI
10.5281/zenodo.21521899.

{[}S5{]} N. Kihara, ``A Thought Experiment on a Model Particle
(Continued) --- The Fifth Frequency and the Composite Curvature Radius
Seen from a Zero-Sum Condition,'' 2026-06-04. Concept DOI
10.5281/zenodo.20535894.

{[}S6{]} N. Kihara, ``A Thought Experiment on a Model Particle --- From
the Phase-Area Floor to the Admissibility Condition for Dimension,''
2026-06-03. Concept DOI 10.5281/zenodo.20528511.

{[}S7{]} N. Kihara, ``Preliminary Summary of Harmonic Readout and c=1
Area Sweep in an AB Two-Body Closed Phase System,'' 2026-07-15. Concept
DOI 10.5281/zenodo.21318696.

{[}S8{]} N. Kihara, ``The Second Thought Experiment: Deriving the
Coulomb-Type Inverse-Square Force from a Product Readout and an Area
Clock - Keeping Two Values and a Linear Law, Obtaining the Focus,
Attraction, Repulsion and Neutrality, and Fixing What Two Values Cannot
Write,'' 2026-09-21. Concept DOI 10.5281/zenodo.22867336.

{[}S9{]} N. Kihara, ``Complex Rotating Wave's Area Cross-Term Curvature
Oscillation Readout: Minimal Application to Gravitational Waves,''
2026-09-01. Concept DOI 10.5281/zenodo.22230940.

{[}S10{]} N. Kihara, ``Paper 0: Distortion of the Geodesic Unit Cell in
Positively-Curved Constant-Curvature Space --- Exact Evaluation of Edge,
Angle, Area, and Volume,'' 2026-07-11. Concept DOI
10.5281/zenodo.20680269.

{[}S11{]} N. Kihara, ``On the Limits of the Curvature Radius in Central
Projection --- The Cases R → ∞ and R → 0,'' 2026-04-12. Concept DOI
10.5281/zenodo.19526548.

{[}S12{]} N. Kihara, ``Regularity of Interspherical Central Projection:
Formulation without Tangent Hyperplanes and Degeneracy at R=0,''
2026-04-22. Concept DOI 10.5281/zenodo.19692191.

{[}S13{]} N. Kihara, ``Self-consistent Determination of the Curvature
Radius and Recovery of the Tangherlini Evaporation Scaling (Paper 5),''
2026-04-28. Concept DOI 10.5281/zenodo.19837595.

{[}S14{]} N. Kihara, ``The Choice of Time Direction and the Stage for
Amplitudes: the Marking Theorem, Derived Z4 Phases, Inversion Holonomy,
and the Canonical Counting Theorem (Dual Geometry of Wavelength and
Frequency Spaces, Paper 14),'' 2026-06-14. Concept DOI
10.5281/zenodo.20665661.

{[}S15{]} N. Kihara, ``Position Had Never Been Read: Sign Sectors,
Duality Theorems, the Readout Hierarchy, and the Derived Clock (Dual
Geometry of Wavelength and Frequency Spaces, Paper 13),'' 2026-06-14.
Concept DOI 10.5281/zenodo.20665633.

{[}S16{]} N. Kihara, ``Conserved Quantities and Uncertainty as
Phase-Space Area: A Unified Viewpoint of Symplectic Symmetry, Wick
Rotation, and Stone's Theorem,'' 2026-06-03. Concept DOI
10.5281/zenodo.20521566.

{[}S17{]} N. Kihara, ``The Twelfth Thought Experiment: From the
Fine-Structure Constant to Observation Maps and Gauge Symmetry -
Starting from the Fine-Structure Constant, Rereading Spacetime, Mass,
Charge, Color Charge and Gauge Symmetry from the Light-Cone Equation and
the Choice of Observable Axes,'' 2026-10-08. Concept DOI
10.5281/zenodo.23226258.

{[}S18{]} N. Kihara, ``Reading Out Time, Space, Mass, Charge,
Translational Velocity and Spin from the Relations and Interference of
Complex Waves without a Background Spacetime: An Integrated Hypothesis
in the Order of the Thought Experiment, from Residents of a Closed
Circle to Residents of a Four-Dimensional Spherical Shell, and on to α,
the Four-Dimensional Lattice and E8,'' 2026-10-15. Concept DOI
10.5281/zenodo.23384928.

{[}S19{]} N. Kihara, ``Emergence of Born Statistics from the
Equidistribution of the Observer-System Beat: A Conjecture on the
Localized-Kernel Model,'' 2026-06-28. Concept DOI
10.5281/zenodo.20967081.

{[}S20{]} N. Kihara, ``Deriving the Form of the Born Distribution from
the Reproducing-Kernel Property of a Localized Odd-Harmonic Wave:
Reducing the Remaining Postulates to the Squaring Rule and Randomness,''
2026-06-27. Concept DOI 10.5281/zenodo.20965526.

{[}S21{]} N. Kihara, ``A Survey of Related Prior Work for the
Observer-System Beat Conjecture: Relation to Prior Programs and the
Locus of the Testable Regime (Supplement),'' 2026-06-28. Concept DOI
10.5281/zenodo.20985749.

{[}S22{]} N. Kihara, ``Discovery of Finite-Order Resonance in Iterated
Exchange Scattering: Identifying Sharp Peaks near
Fine-Structure-Constant Inverse Values 137 and 128 with a Reproducible
Wave-Packet Model,'' 2026-07-18. Concept DOI 10.5281/zenodo.21421366.

{[}S23{]} N. Kihara, ``Lock Dynamics in an Anonymous Two-Channel Closed
Wave System,'' 2026-08-02. Concept DOI 10.5281/zenodo.21763999.

{[}S24{]} N. Kihara, ``Post-Lock Geometry: Double Rotation and the
4-Dimensional Span, and that Plane Switching Is Not an Aggregation
Artifact of Long-Window SVD,'' 2026-09-12. Concept DOI
10.5281/zenodo.22729107.

{[}S25{]} N. Kihara, ``Supplement to the Sixth Thought Experiment: Does
the State Generation Change When the Normalization Reference Is Changed
from Mass to a Charge Unit? - A Restricted Renormalization from G = M =
c = 1 to G = q\_0 = c = 1, and an Exact-Match Verification of the
Initial States of the Five Stored Conditions,'' 2026-09-27. Concept DOI
10.5281/zenodo.22985299.

{[}S26{]} N. Kihara, ``Linear Simplexification by Square-Quantity
Readout and a Dimension-wise Organization of Curvature-Correction
Candidates,'' 2026-06-21. Concept DOI 10.5281/zenodo.20785539.

{[}S27{]} N. Kihara, ``A Map from a Set of Complex Numbers to Spacetime
(1+3): Readout by Quaternions and Partial Projection,'' 2026-10-18.
Concept DOI 10.5281/zenodo.23451494.

{[}S28{]} N. Kihara, ``Zero Closure Was Four-Dimensional: An
Interpretive Note on the Geometry and Algebra of Sum x\_n\^{}2 = 0,''
2026-08-12. Concept DOI 10.5281/zenodo.21902805.

{[}S29{]} N. Kihara, ``Geometric Formulation of 4-Dimensional Space via
Central Projection,'' 2026-04-05. Concept DOI 10.5281/zenodo.19427779.

{[}S30{]} N. Kihara, ``Relational Readout of Quaternions and
Crystallization as Inflation: Anonymity-Preserving Autonomous
Structuring and Partial Projection,'' 2026-10-21. Concept DOI
10.5281/zenodo.23523237.

{[}S31{]} N. Kihara, ``Emergent Condensate, Center of Mass, and
Mass-Like Measure from an Anonymous Complex Standing-Wave Vacuum: An
Anonymous Fourier-Mellin Readout of Condensation Strength, Finite
Extent, and Vacuum-State Gauge Distortion,'' 2026-09-03. Concept DOI
10.5281/zenodo.22282689.

{[}S32{]} N. Kihara, ``Genesis of the Three Spatial Axes and Proper
Time: Re-Derivation of Readouts from the Condensate, and the Uniquely
Unreadable Coordinate Time,'' 2026-08-06. Concept DOI
10.5281/zenodo.21816651.

{[}S33{]} N. Kihara, ``Space, Matter and the Clock Do Not Share a Common
Birth Condition --- What a Seed Must Be to Make a Particle, and the
Lower Bound of Resolution,'' 2026-08-10. Concept DOI
10.5281/zenodo.21874481.

{[}S34{]} N. Kihara, ``Considerations on Mass Structure: A Framework for
Mass Analysis via Axis Scale Values and Signed Areas,'' 2026-04-26.
Concept DOI 10.5281/zenodo.19731609.

{[}S35{]} N. Kihara, ``Formulation of Signed Area: Derivation of Charge
Structure and Consequences of Spin-2 Configurations,'' 2026-04-26.
Concept DOI 10.5281/zenodo.19731599.

{[}S36{]} N. Kihara, ``Mass Structure of Gauge Bosons: Higgs
Non-Involvement Indicated by Sign Vectors,'' 2026-04-24. Concept DOI
10.5281/zenodo.19731607.

{[}S37{]} N. Kihara, ``Geometric Structure of Generation Mixing:
Qualitative Derivation of CKM and PMNS Matrices,'' 2026-04-26. Concept
DOI 10.5281/zenodo.19731601.

{[}S38{]} N. Kihara, ``A Geometric Identity for the Fine-Structure
Constant: From the 4D Unit Ball Volume and its Cube-Packing Deficit
(Paper 7, v3),'' 2026-04-29. Concept DOI 10.5281/zenodo.19869266.

{[}S39{]} N. Kihara, ``Chain Complex Structure on the 4D Hypercubic
Lattice: Structural Correspondence between Kihara Cube-Packing and
Wilson Lattice Gauge Theory (Paper 8, v2),'' 2026-04-29. Concept DOI
10.5281/zenodo.19880467.

{[}S40{]} N. Kihara, ``Geometric Origin of the 0.036 Drift in the α
Self-Consistency Equation: Physical Interpretation via 2D Surface
Vibration Modes in a 4D Hyperball,'' 2026-05-21. Concept DOI
10.5281/zenodo.20319436.

{[}S41{]} N. Kihara, ``Selection of the State-Exchange Weight G\_R = 1 -
R and Candidate Correspondence with the Fine-Structure Constant: A
Numerical Experiment,'' 2026-07-16. Concept DOI 10.5281/zenodo.21396760.

{[}S42{]} N. Kihara, ``Radial Projection: Definition and Relation to
Central Projection,'' 2026-06-06. Concept DOI 10.5281/zenodo.20462569.

{[}S43{]} N. Kihara, ``The Equation of Motion of a Single Wave in a
Relational-Wave Closed System: Kuramoto-Type Phase Coupling on L(K\_N)
and the Condition for a Self-Consistent Floor on the 90-Degree
Phase-Difference Skeleton,'' 2026-09-12. Concept DOI
10.5281/zenodo.22728760.

{[}S44{]} N. Kihara, ``The Inflation Curve Is the Exponential
Amplification of the Linearly Unstable Mode of the Floor: The Largest
Eigenvalue of the Floor Jacobian and Non-Normal, Time-Varying
Corrections,'' 2026-09-12. Concept DOI 10.5281/zenodo.22729095.

{[}S45{]} N. Kihara, ``Spontaneous Symmetry Breaking and Degenerate
Vacua from the 90-Degree Floor: Symmetry Breaking by an Infinitesimal
Seed and the Delta=-2pi/N Lock (Exhaustive N=6,\ldots,40),'' 2026-09-12.
Concept DOI 10.5281/zenodo.22729088.

{[}S46{]} N. Kihara, ``Geometric Rapid Expansion Is Specific to Unstable
Self-Consistent Closures: Causal Discrimination of Onset Modes via
Generic Zero-Closure Initial States,'' 2026-08-05. Concept DOI
10.5281/zenodo.21798854.

{[}S47{]} N. Kihara, ``Mechanism of Inflation-like Rapid Expansion in
Self-Consistent Closed Relational-Wave Systems --- Second Edition:
Correction of Computational Conditions and Re-examination for N =
3--16,'' 2026-08-30. Concept DOI 10.5281/zenodo.22112008.

{[}S48{]} N. Kihara, ``The Eleventh Thought Experiment: Emission Records
of Two-Body Systems without a Background Spacetime - Computing Standard
Theory (Dirac-Coulomb Bound States, 1PN Two-Body Gravity, Einstein A
Coefficients) Referenced Only to the Rest Frame of an Absorber at
Infinity, and Reading Out Emission-Absorption Relations for the Hydrogen
Atom, Electron-Electron, Proton-Proton and the Hydrogen Molecule,''
2026-10-07. Concept DOI 10.5281/zenodo.23199299.

{[}S49{]} N. Kihara, ``The Thirteenth Thought Experiment: Why Does a
Particle That Has Acquired Mass Not Dissipate and Vanish? - Coherence as
a Relation, Harmonic Exchange of Localized Waves, Internal Clocks and
Particle Lifetimes Considered from a Single Interaction,'' 2026-10-10.
Concept DOI 10.5281/zenodo.23266250.

{[}S50{]} N. Kihara, ``The Fourteenth Thought Experiment: Why Can
Different Symmetries Never Be Fully Reconciled? - Non-commensurability,
Pull-in and Metastability of Integer-Selecting Symmetries, and Readout
along the Direction of Evolution,'' 2026-10-12. Concept DOI
10.5281/zenodo.23315557.

\subsubsection{External references}\label{external-references}

{[}X1{]} O. C. Stoica, ``3D-Space and the preferred basis cannot
uniquely emerge from the quantum structure,'' arXiv:2102.08620.

{[}X2{]} M. V. Berry, ``Quantal phase factors accompanying adiabatic
changes,'' Proc. R. Soc. Lond. A 392, 45--57 (1984).

{[}X3{]} M. P. do Carmo, \emph{Differential Geometry of Curves and
Surfaces} (Prentice-Hall, 1976).

{[}X4{]} H. Minkowski, ``Raum und Zeit,'' Physikalische Zeitschrift 10,
104--111 (1909).

{[}X5{]} R. Penrose, W. Rindler, \emph{Spinors and Space-Time}, Vol. 1
(Cambridge University Press, 1984).

{[}X6{]} D. Hestenes, ``Real Spinor Fields,'' J. Math. Phys. 8, 798--808
(1967).

{[}X7{]} D. Hestenes, ``The zitterbewegung interpretation of quantum
mechanics,'' Found. Phys. 20, 1213--1232 (1990).

{[}X8{]} D. Hestenes, ``Spacetime physics with geometric algebra,'' Am.
J. Phys. 71, 691--714 (2003).

{[}X9{]} J. Yvon, ``Équations de Dirac-Madelung,'' J. Phys. Radium 1,
18--24 (1940).

{[}X10{]} T. Takabayasi, ``Relativistic Hydrodynamics of the Dirac
Matter,'' Prog. Theor. Phys. Suppl. 4, 2--80 (1957).

{[}X11{]} E. Madelung, ``Quantentheorie in hydrodynamischer Form,'' Z.
Phys. 40, 322--326 (1927).

{[}X12{]} D. Bohm, ``A Suggested Interpretation of the Quantum Theory in
Terms of `Hidden' Variables. I,'' Phys. Rev.~85, 166--179 (1952).

{[}X13{]} A. Khrennikov, ``Hyperbolic quantum mechanics,''
arXiv:quant-ph/0101002 (2001).

{[}X14{]} A. Khrennikov, ``Representation of the contextual statistical
model by hyperbolic amplitudes,'' J. Math. Phys. 46, 062111 (2005).

{[}X15{]} F. Moldoveanu, ``Non viability of hyperbolic quantum mechanics
as a theory of Nature,'' arXiv:1311.6461 (2013).

{[}X16{]} V. V. Kisil, ``Hypercomplex Representations of the Heisenberg
Group and Mechanics,'' Int. J. Theor. Phys. 51, 964--984 (2012).

{[}X17{]} A. Connes, C. Rovelli, ``Von Neumann algebra automorphisms and
time-thermodynamics relation in generally covariant quantum theories,''
Class. Quantum Grav. 11, 2899--2917 (1994).

{[}X18{]} J. B. Hartle, S. W. Hawking, ``Wave function of the
Universe,'' Phys. Rev.~D 28, 2960--2975 (1983).

{[}X19{]} P. A. M. Dirac, ``The quantum theory of the emission and
absorption of radiation,'' Proc. R. Soc. Lond. A 114, 243--265 (1927).

{[}X20{]} L. Susskind, J. Glogower, ``Quantum mechanical phase and time
operator,'' Physics Physique Fizika 1, 49--61 (1964).

{[}X21{]} D. T. Pegg, S. M. Barnett, ``Phase properties of the quantized
single-mode electromagnetic field,'' Phys. Rev.~A 39, 1665--1675 (1989).

{[}X22{]} T. N. Palmer, ``The Invariant Set Postulate: a new geometric
framework for the foundations of quantum theory and the role played by
gravity,'' Proc. R. Soc. A 465, 3165--3185 (2009).

{[}X23{]} T. N. Palmer, ``Discretization of the Bloch sphere, fractal
invariant sets and Bell's theorem,'' Proc. R. Soc. A 476, 20190350
(2020).

{[}X24{]} G. 't Hooft, \emph{The Cellular Automaton Interpretation of
Quantum Mechanics} (Springer, Fundamental Theories of Physics 185,
2016).

{[}X25{]} W. R. Hamilton, ``On Quaternions; or on a new System of
Imaginaries in Algebra,'' Phil. Mag. (3) 25, 10--13 (1844).

{[}X26{]} J. H. Conway, D. A. Smith, \emph{On Quaternions and Octonions:
Their Geometry, Arithmetic, and Symmetry} (A K Peters, Natick, 2003).

{[}X27{]} J. C. Baez, ``The Octonions,'' Bull. Amer. Math. Soc. 39,
145--205 (2002).

{[}X28{]} G. Terzakis, M. Lourakis, D. Ait-Boudaoud, ``Modified
Rodrigues Parameters: An Efficient Representation of Orientation in 3D
Vision and Graphics,'' J. Math. Imaging Vis. 60, 422--442 (2018).

{[}X29{]} S. R. Marandi, V. J. Modi, ``A preferred coordinate system and
the associated orientation representation in attitude dynamics,'' Acta
Astronautica 15, 833--843 (1987).

{[}X30{]} H. Schaub, J. L. Junkins, ``Stereographic orientation
parameters for attitude dynamics: A generalization of the Rodrigues
parameters,'' J. Astronaut. Sci. 44, 1--19 (1996).

{[}X31{]} P. Kustaanheimo, E. Stiefel, ``Perturbation theory of Kepler
motion based on spinor regularization,'' J. reine angew. Math. 218,
204--219 (1965).

{[}X32{]} H. Hopf, ``Über die Abbildungen der dreidimensionalen Sphäre
auf die Kugelfläche,'' Math. Ann. 104, 637--665 (1931).

{[}X33{]} L. B. Szabados, ``Three-Space from Quantum Mechanics,'' Found.
Phys. 52, 102 (2022).

{[}X34{]} L. B. Szabados, ``Minkowski Space from Quantum Mechanics,''
Found. Phys. 54, 24 (2024).

{[}X35{]} L. B. Szabados, ``Curved spacetimes from quantum mechanics,''
arXiv:2502.07668 (2025).

{[}X36{]} R. D. Schafer, \emph{An Introduction to Nonassociative
Algebras} (Academic Press, 1966; Dover, 1995).

{[}X37{]} A. Hurwitz, ``Ueber die Composition der quadratischen Formen
von beliebig vielen Variablen,'' Nachr. Ges. Wiss. Göttingen,
Math.-Phys. Kl. (1898) 309--316.

{[}X38{]} N. Furey, ``An Algebraic Roadmap of Particle Theories, Part I:
General Construction,'' Ann. Phys. (Berlin) 537, 2400322 (2025).

{[}X39{]} N. Furey, ``An Algebraic Roadmap of Particle Theories, Part
III: Intersections,'' Ann. Phys. (Berlin) 537, 2400324 (2025).

{[}X40{]} T. P. Singh, ``Quantum theory without classical time:
Octonions, and a theoretical derivation of the fine structure constant
1/137,'' Int. J. Mod. Phys. D 30, 2142010 (2021).

{[}X41{]} F. J. Dyson, ``The Threefold Way. Algebraic Structure of
Symmetry Groups and Ensembles in Quantum Mechanics,'' J. Math. Phys. 3,
1199--1215 (1962).

{[}X42{]} Y. Kuramoto, \emph{Chemical Oscillations, Waves, and
Turbulence} (Springer, Berlin, 1984).

{[}X43{]} M. Thümler, S. G. M. Srinivas, M. Schröder, M. Timme,
``Synchrony for Weak Coupling in the Complexified Kuramoto Model,''
Phys. Rev.~Lett. 130, 187201 (2023).

{[}X44{]} S. Lee, L. Braun, F. Bönisch, M. Schröder, M. Thümler, M.
Timme, ``Complexified synchrony,'' Chaos 34, 053141 (2024).

\begin{center}\rule{0.5\linewidth}{0.5pt}\end{center}

\subsection{Appendix A　Reproducing the figures and the
identities}\label{appendix-a-reproducing-the-figures-and-the-identities}

Figures 1--6 were generated by \texttt{図表/\allowbreak{}make\_\allowbreak{}figures.py} in the
same folder. \texttt{python3\ make\_\allowbreak{}figures.py} outputs the figures with
Japanese labels, and \texttt{python3\ make\_\allowbreak{}figures.py\ en} those with
English labels (\texttt{fig*\_\allowbreak{}en}), to \texttt{図表/\allowbreak{}figures/\allowbreak{}} as PNG,
SVG and PDF. The computation and the drawing are the same; only the text
differs. Dependencies: numpy and matplotlib. No random numbers are used.

The identities used in the text are checked by
\texttt{図表/\allowbreak{}verify\_\allowbreak{}identities.py}. C0: \textbar cosh(w/2)\textbar² ＝
\textbar cos(ω/2)\textbar²; C1: sum ＝ cosh φ; C2: difference ＝ cos θ;
C3: the two expressions for \textbar1 ＋ u\textbar²; C4: k ＝
−tanh(w/2); C5: the formula for \textbar k\textbar² and
\textbar k\textbar{} ＝ 1 ⟺ cos θ ＝ 0; C6: agreement with the modified
Rodrigues parameter; C7: the commutator; C8: e\^{}\{ma＋nθ\} ≠
e\textsuperscript{\{ma\}e}\{nθ\}; C9: in complex numbers a and θ enter
only as a sum; C10: tan ＝ cot ＝ 1 at α ＝ π/2. The results are in
\texttt{図表/\allowbreak{}results/\allowbreak{}checks.txt} (all 11 OK in the run of 2026-10-23).
\texttt{図表/\allowbreak{}run\_\allowbreak{}all.sh} runs the figures and the checks in order and
regenerates \texttt{SHA256SUMS.txt}.

\subsection{Appendix B　Dialogue
record}\label{appendix-b-dialogue-record}

The content of this paper took shape in the author's dialogues
(2026-10-21 to 23) after the publication of the previous paper
{[}S30{]}. The course is kept, uncorrected, in the previous paper's
folder \texttt{../\allowbreak{}四元数関係の結晶化とインフレーション\_\allowbreak{}20261021/\allowbreak{}}:
\texttt{思考実験メモ2\_\allowbreak{}結晶化の動力学へ\_\allowbreak{}対称性拘束と複素化蔵本模型\_\allowbreak{}20261021.md}
(§23--§30), the full record of the dialogue with ChatGPT
\texttt{思考実験\_\allowbreak{}対話全記録\_\allowbreak{}無名性と対称性から時空の読み出し\_\allowbreak{}20261022-23.md},
and the record of statements with Claude
\texttt{../\allowbreak{}蔵本引き込みと割り算の罠\_\allowbreak{}思考実験\_\allowbreak{}20261023/\allowbreak{}思考実験\_\allowbreak{}蔵本引き込みと割り算の罠\_\allowbreak{}発言録\_\allowbreak{}20261023.md}
(all in Japanese). The confusion between the notations cosh(w/2) and
cos(ω/2) (the notation of §2.2) is also recorded there.

\subsection{Version history}\label{version-history}

\begin{itemize}
\tightlist
\item
  v1 (2026-10-23): First version. Published on Zenodo (Version DOI
  10.5281/zenodo.23585437, Concept DOI 10.5281/zenodo.23585436).
\end{itemize}

\end{document}
